El objetivo principal del proyecto es diseñar y desarrollar una aplicación web que permita centralizar y digitalizar los principales procesos asociados a la actividad comercial de un distribuidor de productos capilares profesionales, facilitando la gestión de pedidos, el seguimiento de la actividad comercial y la relación diaria con los salones clientes.

A partir de este objetivo general, se plantean los siguientes objetivos específicos:

\begin{enumerate}
    \item \textbf{Centralizar la información operativa del distribuidor:} reunir en una única plataforma los datos de usuarios, clientes profesionales, comerciales, productos, pedidos, visitas, rutas, comunicaciones y operaciones administrativas relevantes.

    \item \textbf{Digitalizar la relación con los clientes profesionales:} ofrecer a los salones un espacio propio desde el que puedan consultar información comercial, acceder al catálogo, revisar sus pedidos y mantener una comunicación más directa con el distribuidor.

    \item \textbf{Facilitar el trabajo diario de los comerciales:} proporcionar herramientas para consultar clientes asignados, planificar visitas, organizar rutas, preparar pedidos, revisar repartos y registrar información relevante durante la actividad comercial.

    \item \textbf{Estructurar un catálogo técnico-comercial:} permitir la gestión de productos, categorías, líneas, subcategorías, recursos asociados y referencias de coloración, de forma que el catálogo pueda servir tanto para la venta como para el asesoramiento técnico.

    \item \textbf{Gestionar el ciclo operativo de pedidos:} cubrir la preparación, confirmación, reparto, entrega y seguimiento de pedidos, incorporando información de cobros y trazabilidad suficiente para los distintos perfiles implicados.

    \item \textbf{Incorporar soporte al seguimiento técnico del salón:} registrar servicios, productos utilizados, tonalidades, resultados y plantillas de trabajo que ayuden a conservar el historial técnico de cada cliente profesional.

    \item \textbf{Mejorar la comunicación comercial y la fidelización:} contemplar promociones, descuentos, formaciones, avisos internos, recordatorios y canales de comunicación que refuercen la relación entre el distribuidor, los comerciales y los clientes.

    \item \textbf{Garantizar seguridad, trazabilidad y separación por roles:} definir un sistema de acceso que diferencie responsabilidades, proteja las operaciones sensibles y registre eventos administrativos relevantes.

    \item \textbf{Preparar una base técnica extensible:} construir el sistema sobre una arquitectura web mantenible, con persistencia relacional, soporte para dispositivos móviles e integración progresiva con servicios externos de apoyo.
\end{enumerate}

Estos objetivos iniciales delimitan la dirección del proyecto y sirven como base para concretar, en el apartado siguiente, el alcance funcional, organizativo y técnico de la aplicación.
